%% =====================================================================
%% REVISED MANUSCRIPT — generated by src/revision/make_manuscript.py
%% All numbers are filled from outputs/revision/stats/numbers.json so that
%% abstract, tables and text cannot diverge. Do not edit numbers by hand;
%% edit the template and regenerate.
%% Class: wlscirep (Scientific Reports). For another journal swap the
%% \documentclass line and the bibliography style; the body is class-agnostic.
%% =====================================================================
\documentclass[fleqn,10pt]{wlscirep}
\usepackage[utf8]{inputenc}
\usepackage[T1]{fontenc}
\usepackage{float}
\usepackage{booktabs}
\usepackage{algorithm}
\usepackage{algpseudocode}
\usepackage{multirow}
\graphicspath{{./}{./figures/}}
\newcommand{\pending}[1]{\textcolor{red}{\textbf{[PENDING: #1]}}}

\title{Leakage, Pseudoreplication, and Model Complexity in Slice-Based Nucleus Segmentation of Confocal Stacks}

\author[1,*]{Varun Kumar Dasoju}
\author[2]{Qingsu Cheng}
\author[1]{Zeyun Yu}
\affil[1]{University of Wisconsin Milwaukee, Department of Computer Science, Milwaukee, WI 53201, USA}
\affil[2]{University of Wisconsin Milwaukee, Department of Biomedical Engineering, Milwaukee, WI 53201, USA}
\affil[*]{vdasoju@uwm.edu}

\keywords{Nucleus segmentation, Confocal microscopy, Mammary epithelial cells, Gabor filters, Class imbalance, Data leakage, UNet++}

\begin{abstract}
Evaluation of segmentation on adjacent confocal sections is vulnerable to both train--test leakage and pseudoreplication. We audited @@N_TOTAL@@ annotated mammary epithelial cell images organized conservatively into @@N_VOL@@ shape-defined source groups. In the original slice-level split, @@LEAK_ORIG_95@@\% of test images had a training image with normalized cross-correlation above 0.95. We therefore separated source groups across partitions, standardized training budgets, averaged metrics within source group, bootstrapped groups, and used paired group-level tests. On @@N_TEST@@ test slices in @@V_TEST@@ held-out groups, the full EfficientNet-B7 UNet++ configuration attained macro Dice @@PROP_DICE@@ (95\% group-bootstrap CI @@PROP_DICE_LO@@--@@PROP_DICE_HI@@); a vanilla EfficientNet-B7 UNet++ attained @@D:vanilla_effb7_unetpp@@, standard baselines ranged from @@WORST_BASE_DICE@@ to @@BEST_BASE_DICE@@, and zero-shot Cellpose-SAM attained @@CELLPOSE_ZEROSHOT_DICE@@. No comparison with the full configuration was significant after Holm correction. Four-fold leave-group-out cross-validation gave macro Dice @@CV_POOLED@@ across all @@N_VOL@@ groups. Gabor input, SCSE attention, complexity-weighted sampling and the stabilized loss produced small, seed-dependent changes rather than consistent gains; encoder performance was also non-monotonic with parameter count. These results show that acquisition-aware grouping, cluster-aware inference and equal optimization budgets can matter more than architectural additions in small slice-based microscopy datasets.
\end{abstract}

\begin{document}
\flushbottom
\maketitle
\thispagestyle{empty}

%% =====================================================================
\section*{Introduction}

Three-dimensional culture of mammary epithelial cells is a standard model for studying how matrix stiffness, radiation and other perturbations alter tissue organization. Quantitative analysis commonly begins with nucleus foreground segmentation in laser-scanning confocal stacks. Because manual delineation is labor-intensive, a practical question is how far a few hundred annotated sections can be stretched and how evaluation should be designed so reported accuracy transfers beyond the images used for training.

An initial analysis reported Dice 0.951 for a multi-component architecture and much lower scores for standard baselines. Re-audit identified two defects with broad relevance to slice-based microscopy. First, random splitting of 2D sections placed highly correlated images across training and test partitions; @@LEAK_ORIG_95@@\% of test images had a training near-duplicate (Fig.~\ref{fig:leak}). Second, a threshold intended for 0/255 masks treated 0/1 masks as empty, corrupting dataset summaries and sampling weights. We rebuilt the evaluation around conservative source groups and treat groups, rather than sections, as the independent statistical units.

We contribute (i) a reproducible correlation audit for detecting leakage among related sections; (ii) a budget-matched comparison of conventional encoder--decoders, Cellpose-SAM and four proposed additions; (iii) source-group bootstrap intervals and paired group-level tests that avoid treating adjacent sections as independent replicates; and (iv) an explicit account of how mask encoding, image loading, sampling and filter normalization affected the initial pipeline. The aim is not to promote another architecture, but to identify which conclusions survive a more defensible evaluation.

\subsection*{Related work}
\textbf{Handcrafted texture features in breast image analysis.} Before end-to-end learning, histopathological breast image classification relied on engineered descriptors---grey-level co-occurrence, local binary patterns, run-length and fractal texture measures---whose comparative value was systematically studied\cite{ozturk2018glcm}, and early CNNs such as HIC-net\cite{ozturk2019hicnet} established that learned features could supersede them. Our Gabor channel sits deliberately between these traditions: an engineered orientation-selective texture descriptor supplied as an input to a learned model, so that its contribution can be measured by ablation.

\textbf{Encoder--decoder segmentation.} U-Net\cite{ronneberger2015unet} and its nested variant UNet++\cite{zhou2018unet++} remain the reference designs; attention gating\cite{oktay2018attention}, SCSE recalibration\cite{roy2018concurrent,hu2018squeeze}, atrous encoders\cite{chen2018deeplab} and compound-scaled encoders\cite{tan2019efficientnet} are the standard components we compare. Self-configuring pipelines\cite{isensee2021nnu} automate these choices.

\textbf{Nucleus-specialist and foundation models.} StarDist\cite{schmidt2018stardist} and HoVer-Net\cite{graham2019hovernet} encode nuclear shape priors; Cellpose\cite{stringer2021cellpose} and its SAM-based successor Cellpose-SAM\cite{pachitariu2025cellposesam}, together with SAM\cite{kirillov2023sam} and MedSAM\cite{ma2024medsam}, offer strong zero-shot cellular segmentation. We include Cellpose-SAM, zero-shot and fine-tuned, as the specialist reference.

\textbf{Class-imbalance objectives and training stability.} Dice\cite{milletari2016vnet}, generalised Dice\cite{sudre2017dice}, focal\cite{lin2017focal} and Tversky\cite{salehi2017tversky,abraham2019novel} losses address foreground scarcity; OneCycle scheduling\cite{smith2019super}, EMA selection\cite{polyak1992acceleration} and gradient clipping\cite{brock2021high} address optimization stability.

\textbf{Gap.} What is missing from the small-data microscopy literature is not another architecture but a demonstration, on the same data, of how much reported accuracy depends on split protocol and training budget, together with component-level evidence under a fair protocol. That is the gap this paper addresses.

%% =====================================================================
\section*{Results}

\subsection*{Dataset grouping and the corrected evaluation protocol}
The dataset (Methods) consists of @@N_TOTAL@@ manually annotated 2D images. Exact image and mask matching links 334 pairs to eight named confocal stacks; source-stack identifiers are not retained for the other images. We therefore use image geometry to form @@N_VOL@@ conservative shape-defined source groups. These groups prevent images with the same geometry from crossing partitions but must not be interpreted as @@N_VOL@@ verified acquisitions. Two files containing image data in place of labels were excluded. Every remaining image contains annotated foreground; the foreground fraction ranges from @@FG_MIN@@\% to @@FG_MAX@@\% (mean @@FG_MEAN@@\%). The previously reported ``60\% empty images'' resulted from thresholding 0/1-encoded masks at 127.

\begin{figure}[ht]\centering
\includegraphics[width=\linewidth]{fig_leakage_audit.pdf}
\caption{Train/test similarity audit. (\textbf{a}) For every test image, the maximum normalized cross-correlation (NCC) with any training image under the original slice-level split (red) and the conservative source-group split (blue). Under the original split the median was @@LEAK_ORIG_MED@@ and @@LEAK_ORIG_95@@\% of test images exceeded 0.95; under the grouped split the median is @@LEAK_NEW_MED@@ and the maximum @@LEAK_NEW_MAX@@. (\textbf{b}) A representative highly correlated pair assigned across partitions by the original split.}
\label{fig:leak}
\end{figure}

Fig.~\ref{fig:leak} quantifies the leakage in the original split. Table~\ref{tab:released} re-evaluates the originally submitted model with its original preprocessing: on the half of the original test set that has a training near-duplicate it scores @@REL_LEAK@@ Dice, on the other half @@REL_NOLEAK@@. The effect is real but modest; the larger consequence of the slice-level split, as the remainder of this section shows, was that it obscured how close all reasonable architectures are to one another when trained properly.

@@TAB_RELEASED@@

All results below use a source-group split: @@N_TRAIN@@ training slices from @@V_TRAIN@@ groups, @@N_VAL@@ validation slices from @@V_VAL@@ groups, and @@N_TEST@@ test slices from @@V_TEST@@ groups. No shape-defined group appears in more than one partition, and the maximum test-to-train NCC is @@LEAK_NEW_MAX@@.

\subsection*{Budget-matched comparison}
Table~\ref{tab:sota} and Fig.~\ref{fig:sota} compare the proposed configuration against seven architectures, every one trained under the identical protocol of 30 epochs with OneCycleLR, AdamW, batch size 8, the same augmentation and the same model-selection rule (Methods). The two baselines omitted from our earlier submission as ``failed to converge'' (MA-Net; UNet++ with ResNet-50) are included; under the shared protocol they converge. Cellpose-SAM is reported zero-shot and after fine-tuning on annotated volumes disjoint from each test volume.

@@TAB_SOTA@@

\begin{figure}[ht]\centering
\includegraphics[width=0.8\linewidth]{fig_sota.pdf}\\[2pt]
\includegraphics[width=0.8\linewidth]{fig_boxplot.pdf}
\caption{Budget-matched comparison on the source-group test set. Top: source-group macro Dice with 95\% group-bootstrap confidence intervals. Bottom: descriptive per-image Dice distributions; images are not treated as independent replicates.}
\label{fig:sota}
\end{figure}

Under equal budgets every architecture converged, including the two previously reported as failures. The full configuration attained source-group macro Dice @@PROP_DICE@@ (95\% group-bootstrap CI @@PROP_DICE_LO@@--@@PROP_DICE_HI@@), IoU @@PROP_IOU@@, precision @@PROP_P@@ and recall @@PROP_R@@, with HD95 @@PROP_HD95@@\,$\mu$m, ASSD @@PROP_ASSD@@\,$\mu$m and normalized surface Dice at 2\,px of @@PROP_NSD@@. Across @@N_SEEDS@@ random seeds its macro Dice was @@PROP_SEED_MEAN@@$\pm$@@PROP_SEED_SD@@. Four-fold leave-group-out cross-validation over all @@N_VOL@@ groups gave macro Dice @@CV_MEAN@@ (between-group SD @@CV_SD@@; fold macro range @@CV_MIN@@--@@CV_MAX@@; @@CV_ROWS@@). The strongest standard baseline, @@BEST_BASE@@, reached @@BEST_BASE_DICE@@; the weakest reached @@WORST_BASE_DICE@@. The paired group-level comparison with the strongest baseline gave Holm-adjusted $p$=@@BEST_BASE_P@@. Zero-shot Cellpose-SAM reached @@CELLPOSE_ZEROSHOT_DICE@@; author fine-tuning increased precision to @@CELLPOSE_FINETUNED_P@@ while reducing recall (Dice @@CELLPOSE_FINETUNED_DICE@@). Headline numbers exclude test-time augmentation.

Throughout, ``full configuration'' denotes the EfficientNet-B7 UNet++ with Gabor input, SCSE, complexity-weighted sampling and the stabilized objective. The simpler vanilla configuration in Table~\ref{tab:ablation} scored @@D:vanilla_effb7_unetpp@@. In the initial analysis, baselines scored 0.47--0.76 and the full configuration scored 0.95; under the standardized protocol the spread is much narrower. A controlled loader experiment showed that the legacy 16-bit-to-8-bit conversion reduced Dice by only about 0.01, so it cannot explain the former gap. Unequal optimization budgets and cross-partition image similarity are documented contributors, but the surviving records do not identify the full cause of the original baseline failures.

\subsection*{Component ablation}
Table~\ref{tab:ablation} changes one element of the full configuration per ablation row; the final vanilla anchor removes all four additions together. Removing the Gabor channel gives @@D:abl_no_gabor@@ ($\Delta$=@@DELTA:abl_no_gabor@@, $p$=@@P:abl_no_gabor@@); replacing it with Sobel, Scharr, Laplacian-of-Gaussian or Canny channels gives @@D:edge_sobel@@, @@D:edge_scharr@@, @@D:edge_log@@ and @@D:edge_canny@@ respectively. Removing SCSE gives @@D:abl_no_scse@@; uniform instead of complexity-weighted sampling gives @@D:abl_uniform_sampling@@; replacing the stabilized objective with plain Dice+BCE gives @@D:abl_loss_dicebce@@ and with a Tversky ($\alpha$=0.3, $\beta$=0.7) plus morphological-boundary objective gives @@D:abl_loss_tversky_boundary@@; a plain U-Net decoder gives @@D:abl_plain_unet@@. The vanilla EfficientNet-B7 UNet++ anchor---RGB input, no SCSE, uniform sampling, plain Dice+BCE---reaches @@D:vanilla_effb7_unetpp@@ ($\Delta$=@@DELTA:vanilla_effb7_unetpp@@, $p$=@@P:vanilla_effb7_unetpp@@).

@@TAB_ABL@@

\begin{figure}[ht]\centering
\includegraphics[width=0.9\linewidth]{fig_ablation.pdf}
\caption{Component ablation. Each bar changes one element of the full configuration; error bars are 95\% bootstrap intervals over held-out source groups.}
\label{fig:abl}
\end{figure}

\subsection*{Gabor parameter sensitivity}
Table~\ref{tab:gabor} varies one parameter of the filter bank at a time. Halving the orientations to 4 gives @@D:gabor_orient4@@ and doubling to 16 gives @@D:gabor_orient16@@; two scales give @@D:gabor_scale2@@ and five give @@D:gabor_scale5@@; $\sigma$=2 gives @@D:gabor_sigma2@@ and $\sigma$=10 gives @@D:gabor_sigma10@@. Fig.~\ref{fig:gabor} shows the actual kernels and their responses on a held-out slice.

@@TAB_GABOR@@

\begin{figure}[ht]\centering
\includegraphics[width=\linewidth]{fig_gabor_bank.pdf}
\caption{The multi-scale Gabor edge channel, computed by the released code on a held-out test slice. Top: six of the 24 kernels ($\theta$, $\lambda$ as labelled). Bottom: input, two individual filter magnitudes, the maximum over all 24 filters that forms the fourth input channel, a Sobel magnitude for comparison, and the ground truth.}
\label{fig:gabor}
\end{figure}

\subsection*{Encoder ladder}
Table~\ref{tab:ladder} and Fig.~\ref{fig:ladder} hold every other component fixed and vary only the encoder. ResNet-34, ResNet-50 and EfficientNet-B0/B3/B5 reach @@D:ladder_resnet34@@, @@D:ladder_resnet50@@, @@D:ladder_effb0@@, @@D:ladder_effb3@@ and @@D:ladder_effb5@@ respectively, against @@PROP_DICE@@ for EfficientNet-B7 (@@PROP_PARAMS@@\,M parameters). Dice does not increase with parameter count over a tenfold range.

@@TAB_LADDER@@

\begin{figure}[ht]\centering
\includegraphics[width=0.5\linewidth]{fig_encoder_ladder.pdf}
\caption{Dice versus parameter count for the encoder ladder (all other components fixed).}
\label{fig:ladder}
\end{figure}

\subsection*{Cross-validation over all source groups}
The hold-out split evaluates @@V_TEST@@ groups. We additionally ran four-fold leave-source-groups-out cross-validation so every one of the @@N_VOL@@ groups was predicted by a model that did not train on that group (Table~\ref{tab:cv}). Source-group macro Dice was @@CV_POOLED@@ across @@CV_NSLICES@@ images, with fold macro means from @@CV_MIN@@ to @@CV_MAX@@. Precision and recall vary markedly across groups and are associated with the two mask encodings. Because the repository does not preserve annotation-session identifiers for all images, this association is descriptive: it does not establish different annotators, sessions, or causal boundary conventions. The heterogeneity nevertheless shows why uncertainty must be estimated over source groups rather than over images.

@@TAB_CV@@

\subsection*{Training behaviour}
Fig.~\ref{fig:curves} shows training loss and held-out-group validation Dice for the @@N_SEEDS@@ seeds. The best validation epoch for seed 0 was @@PROP_BEST_EP@@ of 30 and a full run took approximately @@PROP_MIN@@ minutes on one RTX~5090. No non-finite loss step was recorded. All images contained positive foreground, so the zero-positive branch was not needed for this dataset.

\begin{figure}[ht]\centering
\includegraphics[width=\linewidth]{fig_training_curves.pdf}
\caption{(\textbf{a}) Training loss and (\textbf{b}) validation Dice on held-out source groups for each seed of the full configuration.}
\label{fig:curves}
\end{figure}


\subsection*{Qualitative behaviour}
Fig.~\ref{fig:qual} shows the proposed model's best, median, lower-quartile and worst test slices. Errors concentrate in two situations: dim, out-of-focus nuclei at the edge of the optical section, which the annotator delineated generously and the model segments conservatively (false-negative rims), and small bright debris adjacent to nuclei (isolated false positives). Neither is specific to the proposed model; the same slices are the hardest for every method in Table~\ref{tab:sota}.

\begin{figure}[ht]\centering
\includegraphics[width=\linewidth]{fig_qualitative.pdf}
\caption{Held-out test slices at the best, median, lower-quartile and worst per-image Dice of the proposed model. Columns: input, ground truth, predicted probability, thresholded prediction, and error map (green true positive, red false positive, blue false negative).}
\label{fig:qual}
\end{figure}

%% =====================================================================
\section*{Discussion}

The central finding is that conclusions about small architectural additions changed after the evaluation protocol was corrected. A slice-level split populated the test set with highly correlated training images; the full configuration received more training epochs than its baselines; and defects in mask thresholding and Gabor normalization meant that two components did not operate as described. The legacy image loader cost only about 0.01 Dice in a controlled comparison. Under conservative grouping and source-group inference, the vanilla EfficientNet-B7 UNet++ reached @@D:vanilla_effb7_unetpp@@ macro Dice and the full configuration reached @@PROP_DICE@@ (seed SD @@PROP_SEED_SD@@). None of the component comparisons was significant after correction for multiple testing. The appropriate conclusion is therefore not that an addition is beneficial or detrimental, but that its effect is unresolved and small relative to between-group heterogeneity.

Encoder performance was non-monotonic with size (Table~\ref{tab:ladder}): EfficientNet-B0 with 6.6\,M parameters reached @@D:ladder_effb0@@ and EfficientNet-B7 with 68\,M reached @@PROP_DICE@@, while intermediate encoders occupied the same narrow range. This does not test ImageNet pretraining against random initialization, so we draw no causal conclusion about pretraining. The budget-matched baselines span @@WORST_BASE_DICE@@--@@BEST_BASE_DICE@@, and no paired group-level comparison with the full configuration survives multiplicity correction. The former U-Net (ResNet-50) score of 0.466 becomes @@D:base_unet_r50@@, but only about 0.01 of that discrepancy is explained by the controlled loader experiment; the rest cannot be reconstructed confidently from the surviving records.

The Gabor ablation is not significant, and alternative edge maps (Sobel @@D:edge_sobel@@, Scharr @@D:edge_scharr@@, LoG @@D:edge_log@@, Canny @@D:edge_canny@@) and the parameter sweep (Table~\ref{tab:gabor}) do not show a consistent advantage. With only seven held-out source groups, these experiments support a narrow conclusion: no benefit of the added edge channel was demonstrated.

Precision and surface distance complement Dice, but neither can establish annotation uncertainty without repeat annotation. The association between mask encoding and group-level precision/recall warrants investigation, yet it is confounded with missing provenance and must not be interpreted as inter-annotator disagreement. A blinded multi-annotator study with a pre-specified sample is needed before claiming an annotation-imposed performance ceiling.

\textbf{Limitations.} The images are MCF10A cultures from one laboratory and microscope, and no external cohort was annotated. Full source-stack identifiers are available for only 334 of 855 image--mask pairs, so the remaining images are grouped conservatively by geometry. The primary test partition has only seven source groups, limiting power and precision even though it contains 206 images. Ground truth was produced by one trained annotator with expert review and no repeat-annotation study. The outcome is binary 2D foreground segmentation, not instance segmentation or full-volume 3D validation. Nothing here has been clinically validated and no clinical claims are made.

%% =====================================================================
\section*{Methods}

\subsection*{Dataset}
Confocal $z$-stacks of MCF10A mammary epithelial cultures were acquired in-house (Zeiss LSM, 40$\times$/1.1\,W objective, 405/461\,nm, pixel size 0.2506\,$\mu$m, $z$-step 1.0\,$\mu$m, 16-bit). \pending{Add cell source/authentication, culture and matrix conditions, nuclear stain, exact microscope model and acquisition settings.} Individual $z$-planes were exported as full-frame 2D images, and nucleus foreground was manually delineated by a trained annotator and reviewed by Q.C. \pending{Add annotation software, written boundary rule and adjudication procedure.} The dataset contains @@N_TOTAL@@ valid image--mask pairs after exclusion of two files containing image data in place of labels. Masks stored as either 0/255 or 0/1 are binarized as $m>0$.

\textbf{Provenance audit and source-group split.} SHA-256 matching of image arrays and binarized masks links 334 pairs exactly to eight named annotated stacks. The repository does not retain acquisition identifiers for the other pairs. We therefore grouped images conservatively by frame shape, yielding @@N_VOL@@ shape-defined source groups. A group may contain more than one acquisition and is not called a biological replicate or verified field. Groups were assigned to training (@@V_TRAIN@@ groups, @@N_TRAIN@@ slices), validation (@@V_VAL@@, @@N_VAL@@) and test (@@V_TEST@@, @@N_TEST@@) partitions stratified by foreground content. The assignment and exact-match audit are provided as supplementary machine-readable files. \textbf{Leakage audit.} For every test image we computed the maximum NCC of its $128{\times}128$ thumbnail against all training thumbnails under both the original slice-level split and the source-group split (Fig.~\ref{fig:leak}).

\subsection*{Preprocessing}
Each 16-bit image is normalised per image to its 1st--99.5th percentile range and replicated to three channels. No region-of-interest cropping is applied at any stage; training and evaluation use one identical preprocessing path. Images are resized to $512{\times}512$ for the network; predictions are resampled to native resolution before every metric is computed.

\subsection*{Multi-scale Gabor edge channel}
A bank of $8\times3=24$ Gabor kernels
\[
G(x,y;\lambda,\theta,\sigma,\gamma)=\exp\!\left(-\frac{x'^2+\gamma^2y'^2}{2\sigma^2}\right)\cos\!\left(\frac{2\pi x'}{\lambda}\right),\quad
x'=x\cos\theta+y\sin\theta,\; y'=-x\sin\theta+y\cos\theta,
\]
with $\theta\in\{0,22.5^\circ,\ldots,157.5^\circ\}$, spatial frequencies $f\in\{0.05,0.15,0.25\}$\,px$^{-1}$ (i.e.\ $\lambda=1/f\in\{20,6.67,4\}$\,px), $\sigma=5$, $\gamma=0.5$, phase $\psi=0$, kernel size $31\times31$. Each kernel is L1-normalised (the sum normalisation of our earlier code is unstable because $\sum G\to0$ at high frequency). The edge channel is the per-pixel maximum of the 24 response magnitudes, min--max normalised, and is concatenated to the three intensity channels. A learnable projection Conv$_{3\times3}$(4$\to$16)--BN--ReLU--Conv$_{3\times3}$(16$\to$8)--BN--ReLU--Conv$_{1\times1}$(8$\to$3) maps the 4-channel input to 3 channels so that ImageNet-pretrained encoder weights are retained. Alternative channels (Sobel, Scharr, LoG, Canny) replace the Gabor map in the ablation.

\subsection*{Network}
EfficientNet-B7 encoder\cite{tan2019efficientnet} with ImageNet weights; UNet++ decoder\cite{zhou2018unet++} with channels (256,128,64,32,16) and SCSE\cite{roy2018concurrent} attention at every decoder block; one output logit map. Implementation via Segmentation Models PyTorch\cite{smp}. Baselines use the same library with the encoder and decoder stated in Table~\ref{tab:sota}; ``Attention U-Net'' denotes U-Net with SCSE decoder attention.

\subsection*{Objective}
\begin{algorithm}[H]
\caption{Stabilised Dice--BCE objective (the objective used for all ``stabilised'' rows)}
\begin{algorithmic}[1]
\State $\mathbf{z}\gets\mathrm{clamp}(\mathbf{z},-10,10)$; $\;\mathbf{p}\gets\mathrm{clamp}(\sigma(\mathbf{z}),10^{-7},1-10^{-7})$
\For{each sample $i$ in the batch} $\text{dice}_i\gets(2\sum p_i t_i+1)/(\sum p_i+\sum t_i+1)$ \EndFor
\State $\mathcal{L}_{\rm Dice}\gets1-\mathrm{mean}_i(\text{dice}_i)$
\State $r\gets\sum\mathbf{t}/|\mathbf{t}|$
\If{$r>0$} $w_{\rm pos}\gets\mathrm{clamp}((1-r)/r,1,50)$ \Else\ $w_{\rm pos}\gets1$ \Comment{zero-positive safeguard}\EndIf
\State $\mathcal{L}_{\rm BCE}\gets\mathrm{BCEWithLogits}(\mathbf{z},\mathbf{t};w_{\rm pos})$
\State if either term is non-finite, replace it by 1 \Comment{non-finite fallback}
\State \Return $0.7\,\mathcal{L}_{\rm Dice}+0.3\,\mathcal{L}_{\rm BCE}$
\end{algorithmic}
\end{algorithm}
This objective trained the full configuration and shared-objective comparisons. The ``plain Dice+BCE'' ablation uses $0.5\,\mathcal{L}_{\rm Dice}+0.5\,\mathcal{L}_{\rm BCE}$ with no clamping, adaptive weight or fallback. The ``Tversky+boundary'' ablation uses $0.5(0.5\mathcal{L}_{\rm BCE}^{w}+0.5\mathcal{L}_{\rm Dice})+0.3\,\mathcal{L}_{\rm Dice}(\partial\mathbf{p},\partial\mathbf{t})+0.2\,\mathcal{L}_{\rm Tv}$ with $\mathcal{L}_{\rm Tv}=1-TP/(TP+0.3\,FP+0.7\,FN)$ and $\partial\mathbf{m}=(\mathbf{m}\oplus K)-(\mathbf{m}\ominus K)$ the $3\times3$ morphological gradient.

\subsection*{Complexity-weighted sampling}
With nucleus fraction $r$ and Canny edge-pixel fraction $e$ of each binarised training mask, $w=0.3$ if $r\le0.001$, else $w=(2+3r)\cdot s\cdot(1+e)$ with $s=1.5$ if $r<0.05$ else $1$; weights are normalised to mean 1 and images drawn with replacement with probability $w_i/\sum_j w_j$. Because every image contains nuclei, the effect is to over-sample small-nucleus and boundary-complex slices. The ``uniform'' ablation shuffles without weights.

\subsection*{Shared training protocol}
All encoder--decoder configurations use 30 epochs; AdamW (weight decay $10^{-4}$); OneCycleLR with peak learning rate $3\times10^{-4}$ and 10\% warm-up; batch size 8; bfloat16 autocast; gradient clipping at 0.5; and the same geometric and intensity augmentation. Loss ablations differ only in the objective stated above. The checkpoint with the highest validation Dice on held-out source groups is evaluated once on the test groups. Each reported encoder--decoder configuration is run with three pre-specified seeds. Cellpose fine-tuning follows the separate four-fold protocol below. Training used one RTX~5090 GPU and PyTorch 2.11.

\subsection*{Evaluation}
Metrics are computed for each image at native resolution after thresholding the resampled probability map at 0.5: Dice, IoU, precision, recall, 95th-percentile Hausdorff distance (HD95), average symmetric surface distance (ASSD) in $\mu$m at 0.2506\,$\mu$m/px, and normalized surface Dice at a 2-px tolerance (NSD$_2$). For each encoder--decoder configuration, predictions are averaged over the three seeds at the source-group metric level. Each metric is then averaged across groups, giving equal weight to each group. Ninety-five percent confidence intervals use a percentile bootstrap of source groups (100,000 resamples)\cite{efron1987bca}; seed-to-seed SD is reported separately. Comparisons with the full configuration use paired two-sided Wilcoxon signed-rank tests\cite{wilcoxon1945} on the seven source-group Dice means, Holm-corrected\cite{holm1979} across reported hold-out comparisons, with paired rank-biserial correlation as effect size. Per-image plots are descriptive only. Test-time augmentation is excluded from headline results.

\subsection*{Cellpose-SAM reference}
Cellpose 4.2.1.1\cite{pachitariu2025cellposesam} with the \texttt{cpsam\_v2} checkpoint (flow threshold 0.4; cell-probability threshold 0.0) was run independently on each test image. Fine-tuned rows use four models trained in a leave-two-traceable-stacks-out scheme; each traceable test stack is evaluated by the fold that excluded it, and untraceable groups are evaluated with models that did not use their images. Fine-tuning labels are watershed-derived instances from the binary foreground masks, so only semantic foreground metrics are interpreted here.

\subsection*{Use of AI tools}
No figure in this paper is AI-generated; all image panels are produced by the released code from the study data. A large language model was used as a coding and editing assistant; the authors verified all code and text.

%% =====================================================================
\section*{Data availability}
All images, masks and metadata underlying the reported results, together with code, source-group assignments, random seeds, per-image metrics and model weights, will be available at \pending{public repository URL} and archived at \pending{Zenodo DOI} under \pending{data/code licences}. Any restriction must be stated specifically and checked against the selected journal's data policy before submission.

\section*{Acknowledgements}
This work is partially supported by the U.S. Department of Energy grant \#DE-SC0025403.

\section*{Author contributions statement}
V.K.D.\ designed the methodology, implemented the framework and the revised evaluation protocol, conducted all experiments and statistical analyses, and drafted the manuscript. Q.C.\ acquired the data, supervised annotation and review, and contributed biological interpretation. Z.Y.\ supervised the project and revised the manuscript. All authors reviewed and approved the final manuscript.

\section*{Additional information}
\noindent\textbf{Competing interests:} The authors declare no competing interests.

%% =====================================================================
\begin{thebibliography}{99}
\bibitem{ozturk2018glcm} \"Ozt\"urk, \c{S}.\ \& Akdemir, B. Application of feature extraction and classification methods for histopathological image using GLCM, LBP, LBGLCM, GLRLM and SFTA. \emph{Procedia Comput. Sci.} \textbf{132}, 40--46 (2018).
\bibitem{ozturk2019hicnet} \"Ozt\"urk, \c{S}.\ \& Akdemir, B. HIC-net: a deep convolutional neural network model for classification of histopathological breast images. \emph{Comput. Electr. Eng.} \textbf{76}, 299--310 (2019).
\bibitem{ronneberger2015unet} Ronneberger, O., Fischer, P.\ \& Brox, T. U-Net: convolutional networks for biomedical image segmentation. In \emph{MICCAI} 234--241 (2015).
\bibitem{zhou2018unet++} Zhou, Z., Rahman Siddiquee, M.~M., Tajbakhsh, N.\ \& Liang, J. UNet++: a nested U-Net architecture for medical image segmentation. In \emph{DLMIA/ML-CDS} 3--11 (2018).
\bibitem{oktay2018attention} Oktay, O. et al. Attention U-Net: learning where to look for the pancreas. In \emph{MIDL} (2018).
\bibitem{roy2018concurrent} Roy, A.~G., Navab, N.\ \& Wachinger, C. Concurrent spatial and channel `squeeze \& excitation' in fully convolutional networks. In \emph{MICCAI} 421--429 (2018).
\bibitem{hu2018squeeze} Hu, J., Shen, L.\ \& Sun, G. Squeeze-and-excitation networks. In \emph{CVPR} 7132--7141 (2018).
\bibitem{chen2018deeplab} Chen, L.-C., Zhu, Y., Papandreou, G., Schroff, F.\ \& Adam, H. Encoder--decoder with atrous separable convolution for semantic image segmentation. In \emph{ECCV} 801--818 (2018).
\bibitem{tan2019efficientnet} Tan, M.\ \& Le, Q. EfficientNet: rethinking model scaling for convolutional neural networks. In \emph{ICML} 6105--6114 (2019).
\bibitem{isensee2021nnu} Isensee, F., Jaeger, P.~F., Kohl, S.~A.~A., Petersen, J.\ \& Maier-Hein, K.~H. nnU-Net: a self-configuring method for deep learning-based biomedical image segmentation. \emph{Nat. Methods} \textbf{18}, 203--211 (2021).
\bibitem{schmidt2018stardist} Schmidt, U., Weigert, M., Broaddus, C.\ \& Myers, G. Cell detection with star-convex polygons. In \emph{MICCAI} 265--273 (2018).
\bibitem{graham2019hovernet} Graham, S. et al. Hover-Net: simultaneous segmentation and classification of nuclei in multi-tissue histology images. \emph{Med. Image Anal.} \textbf{58}, 101563 (2019).
\bibitem{stringer2021cellpose} Stringer, C., Wang, T., Michaelos, M.\ \& Pachitariu, M. Cellpose: a generalist algorithm for cellular segmentation. \emph{Nat. Methods} \textbf{18}, 100--106 (2021).
\bibitem{pachitariu2025cellposesam} Pachitariu, M., Rariden, M.\ \& Stringer, C. Cellpose-SAM: superhuman generalization for cellular segmentation. \emph{bioRxiv} 2025.04.28.651001 (2025). \url{https://doi.org/10.1101/2025.04.28.651001}.
\bibitem{kirillov2023sam} Kirillov, A. et al. Segment Anything. In \emph{ICCV} 4015--4026 (2023).
\bibitem{ma2024medsam} Ma, J. et al. Segment anything in medical images. \emph{Nat. Commun.} \textbf{15}, 654 (2024).
\bibitem{milletari2016vnet} Milletari, F., Navab, N.\ \& Ahmadi, S.-A. V-Net: fully convolutional neural networks for volumetric medical image segmentation. In \emph{3DV} 565--571 (2016).
\bibitem{sudre2017dice} Sudre, C.~H., Li, W., Vercauteren, T., Ourselin, S.\ \& Cardoso, M.~J. Generalised Dice overlap as a deep learning loss function for highly unbalanced segmentations. In \emph{DLMIA/ML-CDS} 240--248 (2017).
\bibitem{lin2017focal} Lin, T.-Y., Goyal, P., Girshick, R., He, K.\ \& Doll\'ar, P. Focal loss for dense object detection. \emph{IEEE TPAMI} \textbf{42}, 318--327 (2020).
\bibitem{salehi2017tversky} Salehi, S.~S.~M., Erdogmus, D.\ \& Gholipour, A. Tversky loss function for image segmentation using 3D fully convolutional deep networks. In \emph{MLMI} 379--387 (2017).
\bibitem{abraham2019novel} Abraham, N.\ \& Khan, N.~M. A novel focal Tversky loss function with improved attention U-Net for lesion segmentation. In \emph{ISBI} 683--687 (2019).
\bibitem{smith2019super} Smith, L.~N.\ \& Topin, N. Super-convergence: very fast training of neural networks using large learning rates. \emph{Proc. SPIE} \textbf{11006}, 369--386 (2019).
\bibitem{polyak1992acceleration} Polyak, B.~T.\ \& Juditsky, A.~B. Acceleration of stochastic approximation by averaging. \emph{SIAM J. Control Optim.} \textbf{30}, 838--855 (1992).
\bibitem{brock2021high} Brock, A., De, S., Smith, S.~L.\ \& Simonyan, K. High-performance large-scale image recognition without normalization. In \emph{ICML} 1059--1071 (2021).
\bibitem{smp} Iakubovskii, P. Segmentation Models PyTorch. \url{https://github.com/qubvel/segmentation_models.pytorch} (2019).
\bibitem{efron1987bca} Efron, B. Better bootstrap confidence intervals. \emph{J. Am. Stat. Assoc.} \textbf{82}, 171--185 (1987).
\bibitem{wilcoxon1945} Wilcoxon, F. Individual comparisons by ranking methods. \emph{Biometrics Bull.} \textbf{1}, 80--83 (1945).
\bibitem{holm1979} Holm, S. A simple sequentially rejective multiple test procedure. \emph{Scand. J. Stat.} \textbf{6}, 65--70 (1979).
\end{thebibliography}
\end{document}
